\documentclass[11pt]{article}
\usepackage{amsmath, amssymb, amsfonts, mathtools}
%\usepackage{hyperref}
\usepackage[margin=1in]{geometry}

\newcommand{\tr}{\mathrm{tr}}
\newcommand{\ii}{\mathrm{i}}
\newcommand{\dd}{\mathrm{d}}
\newcommand{\eps}{\varepsilon}
\newcommand{\mO}{\mathcal{O}}
\usepackage{
  jheppub
} % for details on the use of the package, please see the JINST-\
\newcommand{\nn}{\nonumber}
\usepackage{braket}
\newcommand{\lb}{\left(}
\newcommand{\rb}{\right)}
\usepackage{amsmath, amssymb}
\newcommand{\NL}[1]{\textcolor{blue}{#1}}

\begin{document}
  \section{A conjecture}
\NL{this is a second test}
  Consider two energy eigenstates of a CFT $\ket{\mO_\alpha}$ and
  $\ket{\mO_\beta}$ corresponding to Virasoro primaries $\mO_{\alpha}$ and
  $\mO_{\beta}$ of dimensions $h_{\alpha}$ and $h_{\beta}$, respectively.
  Consider a subregion of angular size $A=(-\theta,\theta)$ on a spatial circle
  of circumference $2\pi$, with $\theta=\pi x$. We denote by $\rho_{\alpha}$ and
  $\rho_{\beta}$ the reduced states on $A$ in their corresponding vectors.

  The conjecture is that for the mixture $1/2(\rho_{\alpha}+\rho_{\beta})$ we
  have the following asymptotic behavior of relative entropy
  \begin{eqnarray}
    &&S(x)=S(\rho_\alpha|(\rho_\alpha+\rho_\beta)/2)\nn\\
    &&S(x)=A_0 x^{\delta_0}, \qquad S(1-x)=A_1(1-x)^{\delta_1},\nn\\
    &&\delta_1\leq \delta_0\ .
  \end{eqnarray}

  In particular, we will be interested in the special case of massless 2d free
  boson CFT and energy eigenstates created by the action of vertex operator
  states $V_{\alpha}=e^{i\alpha\phi}$ and $V_{\beta}= e^{i\beta\phi}$ at the
  center of radial quantization.

  \section{Relative entropy for mixed excited states}

  \subsection*{Setup and Replica Trick}

  Consider a conformal field theory (CFT) in $1+1$ dimensions on a cylinder
  $\mathcal{C}_{1}$ of circumference $R$. We focus on a spatial subsystem $A$ defined
  by the interval $(-l/2, l/2)$ at a fixed time slice. We simplify our notation
  by taking $A=(-\theta,\theta)$ on a cylinder of circumference $2\pi$ such that
  $\theta=\pi x$.

  Our states of interest are the reduced density matrices corresponding to
  global excited coherent states. These states are generated by the action of vertex
  operators $V_{\alpha}=:e^{i\alpha \phi}:$ on the vacuum at past infinity in imaginary
  time on the Euclidean cylinder $\mathcal{C}_{1}$:
  \begin{equation}
    |\alpha\rangle = \lim_{u\to i\infty}\frac{V_{\alpha}(u)|\Omega\rangle}{\braket{V_{-\alpha}(-u)V_\alpha(u)}^{1/2}}
  \end{equation}
  The reduced density matrix of the state $|\alpha\rangle$ on the subsystem $A$
  is denoted by $\rho_{\alpha}$. A simple way to keep track of normalization of the
  density matrix to defined normalized operators
  \begin{eqnarray}
    \tilde{V}_{\alpha}(u^+)\tilde{V}_{-\alpha}(u^-)=\frac{V_{\alpha}(u^{+})V_{-\alpha}(u^{-})}{\braket{V_{\alpha}(u^+)V_{-\alpha}(u^-)}_{\mathcal{C}_1}}
  \end{eqnarray}

  We aim to compute the relative entropy $S(\rho || \sigma_{y})$ where
  $\rho = \rho_{\alpha}$ and the reference state is the mixture
  $\sigma_{y}= y\rho_{\alpha}+(1-y) \rho_{\beta}$ for $0\leq y\leq 1$. The standard
  replica trick for relative entropy relates this quantity to the analytic continuation
  of trace functionals that break replica symmetry. Specifically, for the mixed state
  case, we utilize the following identity derived from the generalized replica
  trick for normalized density matrices:
  \begin{eqnarray}
    S\left(\rho_\alpha \Big\| y\rho_\alpha+(1-y)\rho_\beta\right)&&= \partial_{n}
    \log \left[ \frac{\text{tr}(\rho_{\alpha}^{n})}{\text{tr}(\rho_{\alpha}(y\rho_{\alpha}+(1-y)\rho_{\beta})^{n-1})}
    \right]_{n\to 1}\nn\\
    && =\partial_n\log\left[
    \frac{\braket{\prod_{i=1}^{n} \tilde{V}_{\alpha}(u^+_i)\tilde{V}_{-\alpha}(u^-_i)}_{\mathcal{C}_n}}{\braket{\tilde{V}_\alpha(u^+_n)\tilde{V}_{-\alpha}(u^-_n)\mO^{(n-1)}_{\alpha|\beta}}_{\mathcal{C}_n}}\right]_{n=1}
  \end{eqnarray}
  where
  \begin{eqnarray}
    &&\mO_{\alpha|\beta}^{(n-1)}=\sum_{p=0}^{n-1}y^{n-1-p}(1-y)^p\sum_{\substack{w\in\{\alpha,\beta\}^{n-1}\\ N_\beta(w)=p}}
    \tilde{V}_{w_1}(u_1^+)\tilde{V}_{-w_1}(u_1^-)\cdots \tilde{V}_{w_{n-1}}(u_{n-1}^+)\tilde{V}_{-w_{n-1}}(u_{n-1}^-)\nn\\
  \end{eqnarray}
  and $w=(w_{1},\dots,w_{n-1})$ is an ordered word and $N_{\beta}(w)$ counts how
  many letters equal $\beta$, and the points are on the $n$-cover of cylinder
  $\mathcal{C}_{n}$ at $\pm i\infty$ on each sheet.

  The conformal map
  \begin{eqnarray}
    z=\lb\frac{\sin((u-\theta)/2)}{\sin((u+\theta)/2)}\rb^{1/n}
  \end{eqnarray}
  sends the interval $A=(-\theta,\theta)$ to $A=(0,\infty)$ unifomizing
  $\mathcal{C}_{n}$ to the complex plane. Under this transformation Virasoro primaries
  transform as
  \begin{eqnarray}
    &&\mO(u,\bar{u})=\lb\frac{\partial z}{\partial u} \rb^{h_\alpha}\lb\frac{\partial
    \bar{z}}{\partial\bar{u}}
    \rb^{\bar{h}_\alpha}\mO(z,\bar{z})\nn\\
    &&V_\alpha(u_k,\bar{u}_k)=\lb\frac{z_{k,n}\Lambda}{n}\rb^{h_\alpha}\lb\frac{\bar{z}_{k,n}\Lambda}{n}\rb^{\bar{h}_\alpha}V_\alpha(z_{k,n},\bar{z}_{k,n}),
  \end{eqnarray}
  where in our example
  \begin{eqnarray}
    z^\pm_{k,n}=e^{\frac{i\pi}{n}(2k\pm x)},\qquad \Lambda=\frac{\sin(\theta)}{\cos(\theta)-\cos(u)}
  \end{eqnarray}
  Therefore, for normalized operators $\tilde{V}$ the contribution of $\Lambda$
  terms cancel and we find
  \begin{eqnarray}
    \tilde{V}_\alpha(u_k,\bar{u}_k)&&=\lb\frac{z_{k,n}}{n z_{1,1}}\rb^{h_\alpha}\lb\frac{\bar{z}_{k,n}}{n
    z_{1,1}}\rb^{\bar{h}_\alpha}\tilde{V}_\alpha(z_{k,n},\bar{z}_{k,n})\nn\\
    &&=\lb\frac{e^{i\pi(2k\pm x)(1-n)/n}}{n }\rb^{h_\alpha}\lb\frac{e^{-i\pi(2k\pm
    x)(1-n)/n}}{n}\rb^{\bar{h}_\alpha}\tilde{V}_\alpha(z_{k,n},\bar{z}_{k,n})
  \end{eqnarray}
  For scalar (spinless) operators in $2d$ we have $h_{\alpha}=\bar{h}_{\alpha}$
  and we find
  \begin{equation}
    \tilde{V}_{\alpha}(u_{k},\bar{u}_{k})=n^{-2h_\alpha}  \tilde{V}_{\alpha}(z_{k,n}
    ,\bar{z}_{k,n})
  \end{equation}
  Further assuming that $h_{\alpha}=h_{-\alpha}$ the relative entropy replica
  trick in terms of the correlator on the complex plane gives
  \begin{eqnarray}
    &&S(\rho_\alpha\|y\rho_\alpha+(1-y)\rho_\beta)=\partial_n\log\mathcal{F}_n\big|_{n\to 1}\nn\\
    &&\mathcal{F}_n\equiv \left[
    \frac{\braket{\prod_{k=1}^{n} \tilde{V}_{\alpha}(z^+_{k,n})\tilde{V}_{-\alpha}(z^-_{k,n})}}{\braket{\tilde{V}_\alpha(z^+_{n,n})\tilde{V}_{-\alpha}(z^-_{n,n})\tilde{\mO}^{(n-1)}_{\alpha|\beta}}}\right]_{n=1}\nn\\
    &&\tilde{\mO}_{\alpha|\beta}^{(n-1)}=\sum_{p=0}^{n-1}y^{n-1-p}(1-y)^p\: n^{4p(h_\alpha-h_\beta)}
    \sum_{\substack{w\in\{\alpha,\beta\}^{n-1}\\ N_\beta(w)=p}} \prod_{k=1}^{n-1}\frac{V_{w_k}(z_{k,n}^{+})V_{-w_k}(z_{k,n}^{-})}{\braket{V_{\omega_k}(z^+_{1,1})V_{-\omega_k}(z^-_{1,1})}}\nn\\
  \end{eqnarray}
  % For operators $h_\alpha=\bar{h}_\alpha$, and $h_\beta=\bar{h}_\beta$, we have
  % \begin{eqnarray}
  %   \tilde{\mO}_{\alpha|\beta}^{(n-1)}=\sum_{p=0}^{n-1}y^{n-1-p}(1-y)^p\lb \frac{\Lambda^2}{n^2}\rb^{2p(h_\beta-h_\alpha)}\sum_{\substack{w\in\{\alpha,\beta\}^{n-1}\\ N_\beta(w)=p}}
  %   \prod_{k=1}^{n-1} \frac{V_{w_k}(z_k^+)V_{-w_k}(z_k^-)}{\braket{V_{w_k}(z^+)V_{-w_k}(z^-)}}
  % \end{eqnarray}

  \subsection*{Small $x$ limit}

  To simplify our notation we define $C_{\omega}=C^{L}_{\omega(-\omega)}$ for
  $\omega=\alpha,\beta$. In the small $x$ limit, we can expand in the OPE
  channel
  \begin{eqnarray}
    \frac{V_{w_k}(z_{k}^{+})V_{-w_k}(z_{k}^{-})}{\braket{V_{w_k}(z^+)V_{-w_k}(z^-)}}=\lb
    \frac{\sin(\pi x)}{\sin(\pi x/n)}\rb^{2h_{\omega_k}}\lb 1+ C_\omega(2\sin(\pi
    x/n))^{\Delta_L}\mO_L(e^{2\pi i k/n})+\cdots\rb\nn
  \end{eqnarray}
  where $\mO_{L}$ is the lightest Virasoro primary that we assume to have
  dimensions $(h_{L},\bar{h}_{L})$, weight $\Delta_{L}=h_{L}+\bar{h}_{L}$ less than
  two and spin $s_{L}=h_{L}-\bar{h}_{L}$ \footnote{Otherwise the lowest order
  term would come from stress tensor.}. Therefore, we find
  \begin{align}
    \label{eq:F_n}\mathcal{F}_{n} & =\frac{\,1+(2\sin(\pi x/n))^{2\Delta_L}f(n)+\cdots\,}{\displaystyle \sum_{p=0}^{\,n-1}y^{\,n-1-p}(1-y)^{p}\left(\frac{n^2\sin(\pi x/n)}{\sin(\pi x)}\right)^{4p(h_\alpha-h_\beta)}\sum_{\substack{w\in\{\alpha,\beta\}^{\,n-1}\\ N_\beta(w)=p}}\left[\,1+(2\sin(\pi x/n))^{2\Delta_L}g_{w}(n)+\cdots\,\right] }
  \end{align}
  where the leading correction $f(n)$ and $g_{p}(n)$ are given by the following sums
  over OPE coefficients and light operator correlators on the replica manifold:
  \begin{align}\label{eq: OPE of two spinful primaries}
    f(n) & =C_{\alpha}^{2}\sum_{k\neq k'=1}^{n}\langle \mO_{L}(e^{2\pi i k/n})\mO_{L}(e^{2\pi i k'/n})\rangle \nn    \\
         & =C_{\alpha}^{2}\frac{n}{2}\sum_{k=1}^{n-1}\frac{e^{-2\pi i s_L(k+k')/n}}{(2\sin(\pi k/n))^{2\Delta_L}}\ .
  \end{align}
  Rewrite the sum above by fixing the separation $m:=k-k' \ (\mathrm{mod}\ n)$. For
  each $m\in\{1,2,\dots,n-1\}$, set $k'=k-m$ (mod $n$). Then
  \[
    e^{-2\pi i s_L (k+k')/n}= e^{-2\pi i s_L (2k-m)/n}= e^{+2\pi i s_L m/n}\,e^{-4\pi
    i s_L k/n},
  \]
  and the denominator depends only on $m$. Hence the sum factorizes:
  \begin{equation}
    \sum_{m=1}^{n-1}\frac{e^{2\pi i s_L m/n}}{\bigl(2\sin\frac{\pi m}{n}\bigr)^{2\Delta_L}}
    \underbrace{\sum_{k=0}^{n-1} e^{-4\pi i s_L k/n}}_{=:G(n)}.
  \end{equation}
  The inner sum $G(n)$ is a geometric series with ratio $q=e^{-4\pi i s_L/n}$:
  \[
    G(n)=\sum_{k=0}^{n-1}q^{k}=
    \begin{cases}
      n, & q=1,     \\
      0, & q\neq 1.
    \end{cases}
  \]
  Therefore $S(n)=0$ for integer $n$ unless $q=1$, i.e.
  \[
    e^{-4\pi i s_L/n}=1 \quad\Longleftrightarrow\quad \frac{4s_{L}}{n}\in\mathbb{Z}
    .
  \]
  Now, since $4s_{L}/n$ is not an integer, the numerator $e^{2\pi i s_L m/n}$ is
  at most a sign. For $s_{k}=0$, we are left with
  \begin{align}
     & S_{2}(n)=\frac{n}{2}\sum_{k=1}^{n-1}(2\sin(\pi k/n))^{-2\Delta_L},\qquad S_{2}(1)=0,
  \end{align}
  whose analytic continuation of $S_{2}(n)$ was discussed in
  \cite{calabrese2010entanglement}
  \begin{align}
    S'_{2}(1)=\frac{\sqrt{\pi}\Gamma(\Delta_{L}+1)}{4\Gamma(\Delta_{L}+3/2)}\ .
  \end{align}
  In a general CFT, there could be several inequivalent choices of the lightest primary
  operator $\mO_{L}$ that contribute at the same order. In this case, the leading
  correction $f(n)$ would be given by the sum over all such operators, that
  following \cite{calabrese2010entanglement} we denote by
  \begin{align}
    f(n) & =\mathcal{N}S_{2}(n)
  \end{align}
  with $\mathcal{N}=2$ for free boson and $\mathcal{N}=1$ for Ising model. For similicity,
  we set $\mathcal{N}=1$ in the remainder of this calculation.

  Similarly, we have
  \begin{align}
    g_{w}(n)= & C_{\alpha}\sum_{k=1}^{n-1}\frac{C_{\omega_k}\mathcal{N}}{(2\sin(\pi k/n))^{2\Delta_L}}+\sum_{k\neq k'=1}^{n-1}\frac{C_{\omega_k}C_{\omega_{k'}}}{(2\sin(\pi |k-k'|/n))^{2\Delta_L}}, % \nn\\
    % &=C_\alpha C_\omega \frac{2 S_2(n)}{n}+\sum_{k\neq k'=1}^{n-1}
    % \frac{C_{\omega_k}C_{\omega_{k'}}}{(2\sin(\pi |k-k'|/n))^{4h_L}}
  \end{align}

  We can simplify the denominator of $\mathcal{F}_{n}$ by noting that
  \begin{align}
    % &\sum_{\substack{w\in\{\alpha,\beta\}^{\,n-1}\\
    %                        N_\beta(w)=p}} 1={n-1 \choose p},\nn\\
     & \sum_{\substack{w\in\{\alpha,\beta\}^{\,n-1}\\ N_\beta(w)=p}}\left[\,1+(2\sin(\pi x/n))^{2\Delta_L}g_{w}(n)+\cdots\,\right]={n-1 \choose p}+(2\sin(\pi x/n))^{2\Delta_L}\bar{g}_{p}(n)+\cdots\,\nn                                   \\
     & \bar{g}_{p}(n) =\sum_{\substack{w\in\{\alpha,\beta\}^{\,n-1}\\ N_\beta(w)=p}}g_{w}(n)=C_{\alpha}\sum_{k=1}^{n-1}\frac{A_{k}}{(2\sin(\pi k/n))^{2\Delta_L}}+\sum_{k\neq k'=1}^{n-1}\frac{B_{kk'}}{(2\sin(\pi |k-k'|/n))^{2\Delta_L}},
  \end{align}
  where we have exchanged the order of the sum over words and the sum over
  replica indices and defined
  \begin{align}
     & A_{k}(p,n)=\sum_{\substack{w\in\{\alpha,\beta\}^{\,n-1}\\ N_\beta(w)=p}}C_{\omega_k}={n-2 \choose p}C_{\alpha}+{n-2 \choose p-1}C_{\beta}\equiv A(p,n)\nn                                                              \\
     & B_{kk'}(p,n)=\sum_{\substack{w\in\{\alpha,\beta\}^{\,n-1}\\ N_\beta(w)=p}}C_{\omega_k}C_{\omega_{k'}}={n-3 \choose p}C_{\alpha}^{2}+2{n-3 \choose p-1}C_{\alpha}C_{\beta}+{n-3 \choose p-2}C_{\beta}^{2}\equiv B(p,n), % \nn\\
    % &=\frac{(n-1-p)(n-2-p)}{(n-1)(n-2)}C_{\alpha}^2+\frac{2p(n-1-p)}{(n-1)(n-2)}C_{\alpha}C_{\beta}+\frac{p(p-1)}{(n-1)(n-2)}C_{\beta}^2\nn
  \end{align}
  where we performed the sums over words for fixed $k\neq k'$ and found that the
  result is independent of $k$ and $k'$. Therefore, we obtain
  \begin{align}
    g_{p}(n) & =C_{\alpha}A(p,n)\sum_{k=1}^{n-1}\frac{1}{(2\sin(\pi k/n))^{2\Delta_L}}+B(p,n)\sum_{k\neq k'=1}^{n-1}\frac{1}{(2\sin(\pi |k-k'|/n))^{2\Delta_L}}\nn      \\
             & =(C_{\alpha}A(p,n)-B(p,n))\sum_{k=1}^{n}\frac{1}{(2\sin(\pi k/n))^{2\Delta_L}}+B(p,n)\sum_{k\neq k'=1}^{n}\frac{1}{(2\sin(\pi |k-k'|/n))^{2\Delta_L}}\nn \\
             & =\frac{2S_{2}(n)}{n}(C_{\alpha}A(p,n)-B(p,n))+B(p,n)S_{2}(n)\nn                                                                                          % \\
    %& =A(p,n)\frac{2S_{2}(n)}{n}+B(p,n)S_{2}(n)
  \end{align}
  and as a result we can write
  \begin{align}
     & \left[\,1+(2\sin(\pi x/n))^{2\Delta_L}g_{w}(n)+\cdots\,\right]={n-1 \choose p}+(2\sin(\pi x/n))^{2\Delta_L}g_{p}(n)+\cdots\,
  \end{align}

  Defining
  \begin{equation}
    D_{n}=\left(\frac{n^{2}\sin(\pi x/n)}{\sin(\pi x)}\right)^{2(h_\alpha-h_\beta)}
    ,\qquad A_{1}=1,
  \end{equation}
  we note the following identities
  \begin{align}
     & Z(n)\equiv \sum_{p=0}^{\,n-1}y^{\,n-1-p}((1-y)D_{n})^{p}{n-1\choose p}= (y+(1-y) D_{n})^{n-1}                                 \\
     & A_{\alpha}(n)\equiv \sum_{p=0}^{\,n-1}y^{\,n-1-p}((1-y)D_{n})^{p}{n-2\choose p}=y(y+(1-y) D_{n})^{n-2}                        \\
     & A_{\beta}(n)\equiv \sum_{p=0}^{\,n-1}y^{\,n-1-p}((1-y)D_{n})^{p}{n-2\choose p-1}=(1-y)D_{n}(y+(1-y) D_{n})^{n-2}              \\
     & B_{\alpha\alpha}(n)\equiv\sum_{p=0}^{\,n-1}y^{\,n-1-p}((1-y)D_{n})^{p}{n-3\choose p}= y^{2}(y+(1-y)D_{n})^{n-3}               \\
     & B_{\alpha\beta}(n)\equiv\sum_{p=0}^{\,n-1}y^{\,n-1-p}((1-y)D_{n})^{p}{n-3\choose p-1}= y(1-y) D_{n}(y+(1-y)D_{n})^{n-3}       \\
     & B_{\beta\beta}(n)\equiv\sum_{p=0}^{\,n-1}y^{\,n-1-p}((1-y)D_{n})^{p}{n-3\choose p-2}= (1-y)^{2}D_{n}^{2}(y+(1-y)D_{n})^{n-3},
  \end{align}
  For simplicity, we define
  \begin{align}
     & X(n)=(y+(1-y)D_{n}), \qquad X(1)=1\nn                 \\
     & Y_{\alpha\beta}(n)=(y C_{\alpha}+(1-y)D_{n}C_{\beta})
  \end{align}
  so that
  \begin{align}
    B(p,n) & =X(n)^{n-3}\lb C_{\alpha}^{2}B_{\alpha\alpha}(n)+2C_{\alpha}C_{\beta}B_{\alpha\beta}(n)+C_{\beta}^{2}B_{\beta\beta}(n)\rb\nn \\
           & =Y_{\alpha\beta}(n)^{2}X(n)^{n-3}\nn                                                                                         \\
    A(p,n) & = C_{\alpha}A_{\alpha}(n)+C_{\beta}A_{\beta}(n)=Y_{\alpha\beta}(n)X(n)^{n-2}
  \end{align}
  Then, in the denominator we have
  \begin{align}
     & X(n)^{n-1}\lb 1+\lb\frac{2\pi x}{n}\rb^{2\Delta_L}\Delta(n)+\cdots \rb
  \end{align}
  where
  \begin{align}
    \Delta(n) & \equiv \frac{2S_{2}(n)}{n}\frac{(C_{\alpha}A(p,n)-B(p,n))}{X(n)^{n-1}}+\frac{B(p,n)}{X(n)^{n-1}}S_{2}(n)\nn                                                          \\
              & =\frac{2S_{2}(n)}{n}\lb \frac{C_{\alpha}Y_{\alpha\beta}(n)}{ X(n)}-\frac{Y_{\alpha\beta}(n)^{2}}{X(n)^{2}}\rb+S_{2}(n)\lb \frac{Y_{\alpha\beta}(n)^{2}}{X(n)^{2}}\rb \\
              & =S_{2}(n)\lb \frac{Y_{\alpha\beta}(n)}{X(n)}\left(\frac{2C_{\alpha}}{n}\right)+\frac{Y_{\alpha\beta}^{2}(n)}{X(n)^{2}}\lb 1-\frac{2}{n}\rb\rb                        %& g_{1}(n)=\frac{2(y C_{\alpha}+(1-y)D_{n}C_{\beta})}{n(y+(1-y) D_{n})}+\frac{(yC_{\alpha}+(1-y)D_{n}C_{\beta})^{2}}{(y+(1-y) D_{n})^{2}}
  \end{align}
  Putting this back into (\ref{eq:F_n}) we find
  \begin{align}
    \log\mathcal{F}_{n}=(1-n)\log X(n)+\log\left[\frac{1+\lb \frac{2\pi x}{n}\rb^{2\Delta_L}S_{2}(n)\lb C_{\alpha}^{2}+\cdots\rb}{1+\lb \frac{2\pi x}{n}\rb^{2\Delta_L}\Delta(n)+\cdots}\right]
  \end{align}
  We first note that taking the $n$-derivative of the first factor in $\mathcal{F}
  _{n}$ and sending $n\to 1$ vanishes
  \begin{align}
    \partial_{n}\left[(1-n)\log(y+(1-y)D_{n})\right]\Big|_{n=1}=0
  \end{align}
  and we are left with
  \begin{align}
    S(\rho_{\alpha}\|y\rho_{\alpha}+(1-y)\rho_{\beta}) & =\partial_{n}\left[\lb \frac{2\pi x}{n}\rb^{2\Delta_L}(S_{2}(n) C_{\alpha}^{2}-\Delta(n))+\cdots\right]_{n=1}
  \end{align}
  Since $S_{2}(1)=0$ and $X(1)=1$, we have
  \begin{align}\label{eq: relative entropy in small x limit}
    S(\rho_{\alpha}\|y\rho_{\alpha}+(1-y)\rho_{\beta}) & = (2\pi x)^{2\Delta_L}S_{2}'(1)(C_{\alpha}^{2}-2C_{\alpha}Y_{\alpha\beta}(1)+Y_{\alpha\beta}(1)^{2})+\cdots\nn                \\
                                                       & =(2\pi x)^{2\Delta_L}\frac{\sqrt{\pi}\Gamma(\Delta_{L}+1)}{4\Gamma(\Delta_{L}+3/2)}(1-y)^{2}(C_{\alpha}-C_{\beta})^{2}+\cdots
  \end{align}
  where in the second line we have used $S_{2}(1)=0$. In the special case,
  $\alpha =\beta$ or $y\to 1$ the term vanishes as expected. 
  
  
  
  -----
  
  
  In the case of vertex
  operators in free massless bosons, we have
  \begin{align}
    V_{\omega}(z,\bar{z})V_{-\omega}(0,0)=|z|^{-2\Delta_\omega}\lb 1+i\omega\lb z J(0)+\bar{z}\bar{J}(0)\rb+\omega^{2} (z\bar{z}J(0)\bar{J}(0)+\cdots)\rb
  \end{align}


  In this OPE expansion, consider the vertex operators sitting at $z_k^+, z_k^-$,
  \begin{align}
      \frac{V_{\omega_k}(z_k^+ , \bar{z}_k^+) V_{-\omega_k}(z_k^- , \bar{z}_k^-)}{ \langle V_{\omega_k}(z^+ , \bar{z}^+ ) V_{-\omega_k}(z^- , \bar{z}^-)\rangle}  &= \lb
    \frac{\sin(\pi x)}{\sin(\pi x/n)}\rb^{2h_{\omega_k}}  
    \Bigl( 1+i \omega_k (z_k^+ - z_k^-) J(e^{2\pi i k/n})  \nonumber \\
    & \hspace{2cm} + i\omega_k (\bar{z}_k^+ - \bar{z}_k^-) \bar{J}(e^{-2\pi i k/n}) + \cdots \Bigr) \nonumber \\
    &=
    \lb
    \frac{\sin(\pi x)}{\sin(\pi x/n)}\rb^{2h_{\omega_k}}
    \Bigl( 1- \omega_k (2\sin(\pi x /n))e^{2\pi i k/n} J(e^{2\pi i k/n})  \nonumber \\
    &\hspace{2cm}+ \omega_k (2\sin(\pi x /n))e^{-2\pi i k/n} \bar{J}(e^{-2\pi i k/n})  +\cdots \Bigr)
  \end{align} 

We use the correlators of currents J: $(h,\bar{h}) = (1,0)$ and $\bar{J}$: $(h,\bar{h}) = (0,1)$, with scaling dimensions 1, and spins 1 and -1 respectvely,
\begin{align}
    \langle J(z) J(w) \rangle &= \frac{1}{(z-w)^2} \nonumber \\
    \langle \bar{J}(\bar{z}) \bar{J}(\bar{w}) \rangle&= \frac{1}{(\bar{z} - \bar{w})^2} \nonumber \\
    \langle J(z) \bar{J}(\bar{w}) \rangle &= 0
\end{align}
Also, we have,
\begin{align}
    (z_k - z_{k'})^2 =  -4 e^{\frac{2 \pi i (k+k')}{n}} \sin^2 \lb \frac{\pi (k-k')}{n} \rb
\end{align}
Then we need to use the correlator of the current operators corresponding to different replica sheets ($k \neq k'$),
\begin{align}
    \langle J_{\omega_k}(z_k^+)J_{\omega_{k'}}(z_{k'}^-) \rangle \sim \frac{1} {(z_k - z_{k'})^2}  \nonumber \\
    \implies  \omega_k \omega_{k'} \lb 2 \sin \lb \frac{\pi x}{n}\rb \rb^2 e^{2 \pi i (k+k')/n} \langle J_{\omega_k}(z_k^+)J_{\omega_{k'}}(z_{k'}^-) \rangle &\sim   \frac{\omega_k \omega_{k'}\lb 2 \sin \lb \frac{\pi x}{n}\rb \rb^2 e^{2 \pi i (k+k')/n}}{(z_k - z_{k'})^2} \nonumber \\ 
    & = - \frac{\omega_k \omega_{k'}\lb 2 \sin \lb \frac{\pi x}{n}\rb \rb^2 } {4 \sin ^2 \lb \frac{\pi (k-k')}{n} \rb }
\end{align}
This looks similar to the previous case, where we considered correlators of scalar primaries with constant OPE coefficients, such as in \ref{eq: OPE of two spinful primaries}.

We have $\mathcal{N} = 2$ in 1+1D free boson CFT, so the final answer is twice the answer given in \ref{eq: relative entropy in small x limit}. 

\subsection*{Small $1-x$ limit}

\begin{align}
\frac{
V_{\omega_k}(z_k^+ , \bar{z}_k^+ )
V_{-\omega_{k+1}}(z_{k+1}^- , \bar{z}_{k+1}^-)
}{
\langle
V_{\omega_k}(z^+ , \bar{z}^+ )
V_{-\omega_{k+1}}(z^- , \bar{z}^-)
\rangle
}
&=
\left(
\frac{\sin(\pi x)}{\sin(\pi x/n)}
\right)^{2h_{\omega_k}}
V_{\omega_k-\omega_{k+1}}
(z_{k+1}^- , \bar{z}_{k+1}^-)
\nonumber\\
&\qquad \times
\Bigl(
1
+i\color{blue}\omega_k\color{black}
(z_k^+ - z_{k+1}^-)
J(e^{2\pi i (k+1)/n})
\nonumber\\
&\qquad\qquad
+i\color{blue}\omega_k\color{black}
(\bar z_k^+ - \bar z_{k+1}^-)
\bar J(e^{-2\pi i (k+1)/n})
+\cdots
\Bigr).
\end{align}



\subsection*{Exact solution}
Final answer for 1+1D free boson CFT,
  \begin{equation}
    \boxed{ S\!\left(\rho_\alpha\,\Big\|\,y\rho_\alpha+(1-y)\rho_\beta\right) = -\log\!\left( y+(1-y)\exp\!\big[-(\alpha-\beta)^2(1-\pi x\cot(\pi x))\big] \right). }
    \label{eq:S-final}
  \end{equation}

Finally, we note that in the small interval limit $x\ll 1$ we have
  \begin{equation}
    K(x)=1-\pi x\cot(\pi x)=\frac{\pi^{2}}{3}x^{2}+\frac{\pi^{4}}{45}x^{4}+O(x^{6}
    ).
  \end{equation}
  Hence
  \begin{equation}
    S(y)= (1-y)\,(\alpha-\beta)^{2}\frac{\pi^{2}}{3}x^{2}+O(x^{4}).
  \end{equation}
  The relative entropy vanishes as $x^{2}$ for small intervals, as expected.  
\end{document}